\section{實作重點}

\subsection{Audio CODEC 控制}

WM8731 需要先透過 I2C 進行暫存器設定，設定內容包含輸入來源、輸出音量、資料格式、取樣頻率與啟用狀態。完成設定後，FPGA 再透過音訊資料介面接收或傳送 16 bit 音訊資料，以達成錄音與播放功能。

\subsection{SRAM 與 FLASH 分工}

SRAM 讀寫速度較快，適合作為錄音過程中的即時暫存記憶體；FLASH 具有非揮發特性，適合用於保存確認後的錄音資料。因此本系統採用先寫入 SRAM，再由使用者確認後寫入 FLASH 的方式，提升系統操作彈性，也避免錄音失敗時直接覆蓋 FLASH 資料。

\subsection{FFT 顯示簡化}

由於完整 FFT 運算較複雜，本專題可將 FFT 結果簡化為 8 個頻帶能量，對應至 8$\times$8 LED 矩陣的 8 個欄位。每個欄位高度代表該頻帶音量大小，使 LED 矩陣能呈現類似音樂頻譜分析器的效果。

%========================
% 十四、實驗結果
%========================
\section{實驗結果}

本章節可放入實際 FPGA 執行結果照片與測試畫面。建議放置以下內容：

\begin{enumerate}
\item 系統開機後 LCD 主選單畫面。
\item 錄音模式下 LCD 顯示 Recording 畫面。
\item SRAM 暫存錄音資料並播放確認的畫面。
\item 使用者確認後，系統將 SRAM 資料寫入 FLASH 的畫面。
\item 從 FLASH 讀取音訊資料並播放的測試畫面。
\item 播放音訊時，8$\times$8 LED 矩陣顯示 FFT 頻譜的畫面。
\item Quartus 編譯成功畫面。
\item ModelSim 模擬波形，例如 FSM 狀態轉移或 SRAM 讀寫時序。
\end{enumerate}

\begin{figure}[H]
\centering
\fbox{\parbox[c][5cm][c]{10cm}{\centering 請放入 FPGA 實作照片}}
\caption{FPGA 系統實作畫面}
\end{figure}

\begin{figure}[H]
\centering
\fbox{\parbox[c][5cm][c]{10cm}{\centering 請放入 LCD 顯示畫面}}
\caption{LCD 顯示畫面}
\end{figure}

\begin{figure}[H]
\centering
\fbox{\parbox[c][5cm][c]{10cm}{\centering 請放入 LED 矩陣 FFT 頻譜畫面}}
\caption{LED 矩陣 FFT 頻譜顯示}
\end{figure}